\section{Experience}

\outerlist{

\entrybig
    {\textbf{University of Illinois Urbana-Champaign}}{Urbana, IL (Onsite, Full-Time)}
    {Graduate Research Assistant, xLab}{Fall 2024 – Present}
    \innerlist{
        \item Co-authored \textbf{STRATUS}, an LLM-based multi-agent system enabling autonomous Site Reliability Engineering (SRE) for cloud services; formalized a novel Transactional Non-Regression (TNR) safety specification and demonstrated 1.5–5.4× improvement over state-of-the-art AI SRE agents on AIOpsLab and ITBench benchmarks.
        \item Lead maintainer of \textbf{AIOpsLab}, an open-source framework for benchmarking LLM-based AIOps agents (700+ GitHub stars).
        \item Developed and deployed microservice-based testbeds for AI-driven incident detection, localization, and mitigation.
        \item Integrated and expanded \textbf{Chaos Engineering} techniques, including eBPF-based fault injection to simulate realistic hardware and software failures.
        \item Co-authored \textbf{ITBench} (with IBM), focusing on benchmarking AI agents in Site Reliability Engineering and cloud automation.
        \item Designed and instrumented cloud-native applications to generate real-time telemetry data (metrics, logs, and traces) for AI-based observability.
        \item Built Kubernetes-based environments with automated workload orchestration and fault injection for testing AI models in cloud reliability.
        \item Applied Large Language Models (LLMs) to SRE workflows, enabling automated fault diagnosis and self-healing cloud systems.
        \item Experienced in infrastructure-as-code tools (Helm, Terraform) and cloud-native monitoring stacks (Prometheus, Jaeger, OpenTelemetry).
    }

\entrybig
    {\textbf{NASA Marshall Space Flight Center}}{Huntsville, AL (Onsite, Full-Time)}
    {Software Engineer Intern, Huntsville Operations Support Center (HOSC)}{Summer 2023}
    \innerlist{
        \item Developed a forward error correction library in C++ based on CCSDS standards (BCH, LDPC, and convolutional codes) for NASA's GPDM mission.
        \item Contributed to software automation tools supporting flight controllers.
    }

\entrybig
    {\textbf{NASA Ames Research Center}}{Mountain View, CA (Remote, Full-Time)}
    {Research Intern, Intelligent Systems Division}{Summer 2023}
    \innerlist{
        \item Led reinforcement learning research for multi-agent pathfinding (MAPF) in urban air traffic management, leveraging Python, PyTorch, TensorFlow, and Docker on a supercomputing cluster.
    }

\entrybig
    {\textbf{NASA Goddard Space Flight Center}}{Greenbelt, MD (Remote, Part-Time)}
    {Software Engineer and Data Science Intern, HEASARC}{Fall 2022 – Spring 2023}
    \innerlist{
        \item \textbf{Spring Project:} Trained ML models (scikit-learn, TensorFlow) to classify neutron star spectra using NICER payload data.
        \item \textbf{Fall Project:} Redesigned SkyView, a React + Flask data catalog system for querying astrophysics datasets by sky position.
    }

\entrybig
    {\textbf{NASA Marshall Space Flight Center}}{Huntsville, AL (Onsite, Full-Time)}
    {Software Engineer Intern, Payload Operations Integration Center}{Summer 2022}
    \innerlist{
        \item Designed an automation tool for ISS flight controllers to streamline timeline change processes.
    }

\entrybig
    {\textbf{Freelance}}{Remote, Illinois}
    {Software Engineer}{Sept 2020 – Sept 2023}
    \innerlist{
        \item Earned over \$60,000 on Upwork developing cross-platform apps using Electron.js, React, and Express; managed full client lifecycle from proposal to delivery.
    }

\entrybig
    {\textbf{RoboThink}}{Remote, Illinois}
    {Software Engineer}{June 2018 – Sept 2020}
    \innerlist{
        \item Lead developer for RoboThink’s educational tools built with Electron.js and React.
    }

}
